% 数据结构课程设计报告模板
% 编译方式：XeLaTeX

\documentclass[a4paper,12pt,twoside]{ctexrep}

% 版面与基本宏包
\usepackage{geometry}   %设置页边距的宏包
\geometry{left=2cm,right=2cm,top=2.5cm,bottom=2.5cm}  %设置 上、左、下、右 页边距

\usepackage{setspace}
\usepackage{graphicx}
\usepackage{amsmath,amssymb}
\usepackage{hyperref}

\usepackage{xcolor}
\usepackage[normalem]{ulem}
\usepackage{array}
\usepackage{fancyhdr}
\usepackage{caption}

\usepackage{listings}
\usepackage[most]{tcolorbox}
\tcbuselibrary{listings,skins,breakable}

\usepackage{booktabs}
\usepackage{enumitem}

\usepackage{tikz}
\usetikzlibrary{shapes.geometric, arrows.meta, positioning, fit, backgrounds, calc}

% =========================
% 列表样式（1ch 视作 1\ccwd）
% =========================

\setlist[itemize]{%
  label=\textbullet,
  labelindent=1\ccwd,   % bullet 起始位置：1ch
  labelwidth=1\ccwd,    % 为 bullet 预留 1ch 宽度
  labelsep=0pt,         % 不额外加缝隙
  leftmargin=2\ccwd,    % 正文起始位置：2ch
  itemindent=0pt,
  align=parleft,        % 悬挂对齐：换行与正文起始对齐
  listparindent=0pt,
  topsep=2pt,itemsep=1pt,parsep=0pt,partopsep=0pt
}

\setlist[enumerate]{%
  label=\arabic*.,
  labelindent=2\ccwd,   % 编号起始位置：2ch
  labelwidth=1\ccwd,    % 为编号预留 1ch
  labelsep=0pt,
  leftmargin=3\ccwd,    % 正文起始位置：3ch
  itemindent=0pt,
  align=parleft,
  listparindent=0pt,
  topsep=2pt,itemsep=1pt,parsep=0pt,partopsep=0pt
}


\usepackage{letltxmacro}
\usepackage[numbers]{natbib}
\usepackage{lipsum}

\usepackage[titles]{tocloft}
\setlength{\cftbeforechapskip}{0pt}
\renewcommand{\cftchapleader}{\cftdotfill{\cftchapdotsep}}
\renewcommand{\cftchapdotsep}{\cftdotsep}
\renewcommand{\cftdotsep}{1.0}

\newcolumntype{R}[1]{>{\raggedleft\arraybackslash}p{#1}}
\newcolumntype{C}[1]{>{\centering\arraybackslash}p{#1}}

\usepackage{ifoddpage}

\newcommand{\evenblank}{%
  \clearpage
  \checkoddpage
  \ifoddpage
    \thispagestyle{empty}\null\newpage
  \fi
}

\newcommand{\OneEm}{\symbol{12288}}

\usepackage{fontspec,xltxtra,xunicode}

\defaultfontfeatures{Mapping=tex-text} %如果没有它，一些 tex 特殊字符无法正常使用，比如连字符
% 文章内中文自动换行，可以自行调节
\XeTeXlinebreaklocale "zh"
\XeTeXlinebreakskip = 0pt plus 0.1em

% 中文字体支持
% 字体需要根据自己电脑系统来设置
\usepackage{xeCJK}

\xeCJKsetup{
  CJKglue   = {\hskip 0pt plus 0.18em}, % 靠 stretch 做两端对齐扩字距
  CJKecglue = {\hskip 0pt plus 0.18em}
}

\setmainfont{Times New Roman}        
\IfFontExistsTF{SimSun}{          % 宋体/常规字体 
  \setCJKmainfont{SimSun}
}{
  \setCJKmainfont{Songti SC}
}
\setsansfont{Arial}
\setCJKsansfont{SimHei}              % 黑体/无衬线体 
\setmonofont{Courier New}
\setCJKmonofont{STFangsong}       % 仿宋/等宽字体 

% 定义用于特定位置的中文字体别名 (可选，模版中已包含)
\setCJKfamilyfont{zhhei}{SimHei}
\newcommand*{\hei}{\CJKfamily{zhhei}}
\setCJKfamilyfont{zhkai}{STKaiti}
\newcommand*{\kai}{\CJKfamily{zhkai}}
\setCJKfamilyfont{enroman}{Times New Roman}
\newcommand*{\mytimes}{\CJKfamily{enroman}}
\setCJKfamilyfont{STFangsong}{STFangsong}


% 标题字体：一级、二级标题用 黑体 + Arial
\newcommand{\HeadingFont}{\heiti\sffamily}

\ctexset{
  chapter = {
    format = \HeadingFont\LARGE\centering,
    name = {第,章},
    number = \chinese{chapter},
    beforeskip = 0.2\baselineskip,
    afterskip = \baselineskip
  },
  section = {
    format = \HeadingFont\Large
  },
  % 三级及以下标题采用 宋体 + Times New Roman
  subsection = {
    format = \heiti\sffamily\large
  },
  subsubsection = {
    format = \songti\bfseries\normalsize
  }
}

% 正文行距
\onehalfspacing

\setlength{\parskip}{0pt}% 段前段后不加额外空白
\setlength{\parindent}{2\ccwd} % 首行缩进 2 字（常见 Word 风格）

% 代码环境设置：仿宋 + Courier，字号约 10.5pt（用 \small 近似）
\newtcblisting{codeblock}{
  listing only,
  listing engine=listings,
  breakable,
  enhanced,
  colback=black!3,
  colframe=black!0,
  boxrule=0pt,
  arc=2pt,
  drop shadow,
  left=6pt,right=6pt,top=6pt,bottom=6pt,
  listing options={
    basicstyle=\ttfamily\small,
    columns=fullflexible,
    keepspaces=true,
    breaklines=true,
    showstringspaces=false
  }
}

% 单行字段容器
\newcommand{\FillLine}[1]{%
  \uline{\hspace{\linewidth}}%
  \llap{#1}%
}

% 多行字段容器
\newcommand{\Field}[2][8cm]{%
  \begin{minipage}[t]{#1}
    \setlength{\parindent}{0pt}% 取消缩进，避免第一行缩进
    \centering #2%
  \end{minipage}%
}

% 居中带下划线的文本字段容器
\newcommand{\CenteredField}[2][8cm]{%
  \begin{minipage}[t]{#1}
    \setlength{\parindent}{0pt}
    \centering
    \underline{\hbox to #1{\hfil#2\hfil}}
  \end{minipage}%
}

% 设置页眉页脚样式
\pagestyle{fancy}
\fancyhf{} % 清除所有页眉页脚的默认设置
\setlength{\headheight}{14pt}

% 页码：装订线外侧（外侧=奇数页右、偶数页左）
\fancyhead[RO,LE]{\thepage}

% 页眉中央：报告标题（组合）
\fancyhead[CE,CO]{\kai \ReportTitleA\ReportTitleB}

\renewcommand{\headrulewidth}{0.4pt} % 添加页眉线
\renewcommand{\footrulewidth}{0pt} % 页脚线保持移除

% 定义一个用于摘要页的样式：只显示页眉，不显示页码（页脚为空）
% 配合 \pagenumbering{gobble} 使用
\fancypagestyle{fancy_no_page}{
  \fancyhf{}
  \fancyhead[RO]{\kai 数据结构课程设计报告} 
  \fancyhead[LE]{\kai \ReportTitleA\ReportTitleB} 
  \renewcommand{\headrulewidth}{0.4pt} % 保持页眉线
  \renewcommand{\footrulewidth}{0pt}   % 保持无页脚线
}

% 章首页
\fancypagestyle{plain}{
  \fancyhf{}
  \fancyhead[RO,LE]{\thepage}
  \fancyhead[CE,CO]{\kai \ReportTitleA\ReportTitleB}
  \renewcommand{\headrulewidth}{0.4pt}
  \renewcommand{\footrulewidth}{0pt}
}

% 设置浮动体位置参数
\setcounter{topnumber}{5}
\setcounter{bottomnumber}{5}
\setcounter{totalnumber}{10}
\renewcommand{\topfraction}{0.95}
\renewcommand{\bottomfraction}{0.95}
\renewcommand{\textfraction}{0.05}

\makeatletter
\setlength{\@fptop}{0pt}
\setlength{\@fpbot}{0pt plus 1fil}
\makeatother

% 自定义图表编号格式: x-y (章节-序号)
\renewcommand{\thefigure}{\arabic{chapter}-\arabic{figure}}
\renewcommand{\thetable}{\arabic{chapter}-\arabic{table}}

% 设置图表标题格式
\DeclareCaptionFormat{myformat}{\songti#1#2 #3}
\captionsetup[figure]{format=myformat,labelfont=rm,labelsep=none,position=below}
\captionsetup[table]{format=myformat,labelfont=rm,labelsep=none,position=top}
\captionsetup[lstlisting]{format=myformat,labelfont=bf,labelsep=none,position=top}

% 设置图表标题字体和大小
\captionsetup{font={small,singlespacing}}
\captionsetup[figure]{name=图}
\captionsetup[table]{name=表}


%%%%%%%%%%%%%%%%%%%%%%%%%%%%%
%%%%%%%%% 填写报告信息 %%%%%%%%
%%%%%%%%%%%%%%%%%%%%%%%%%%%%%

\newcommand{\ReportTitleA}{教室预约系统}
\newcommand{\ReportTitleB}{} % 如果标题过长可在第二行手动续行
\newcommand{\StudentClass}{24 软件四}
\newcommand{\StudentNameA}{步{\OneEm}星辰（24030601）}
\newcommand{\StudentNameB}{王{\OneEm}炀成（24030625）}
\newcommand{\StudentNameC}{} % 如果不需要这行可以留空
\newcommand{\Supervisor}{叶鸿}
\newcommand{\Score}{}
\newcommand{\SubmitDate}{2026 年 1 月 9 日}
\begin{document}

%%%%%%%%%%%%%%%%%%%%
% 封面（不显示页码）
%%%%%%%%%%%%%%%%%%%%

\begin{titlepage}
  \thispagestyle{empty}

  % 上部：Logo
  \vspace*{1cm}
  \begin{center}
    \includegraphics[width=0.7\textwidth]{assets/cit_logo.pdf}
  \end{center}

  % 中部：课程名与报告名
  \vfill

  \begin{center}
    {\HeadingFont\fontsize{36pt}{40pt}\selectfont 数据结构课程设计}
  \end{center}

  % 下部：键值对信息表（无网格线）
  \vfill

  \begin{center}
  {\large
  \renewcommand{\arraystretch}{1.6} % 行距

  \begin{tabular}{@{}R{3cm}C{10cm}@{}} % 列宽
    项目名称： &
      \CenteredField[10cm]{\ReportTitleA} \\
              &
      \CenteredField[10cm]{\ReportTitleB} \\
    班{\OneEm}{\OneEm}级： &
      \CenteredField[10cm]{\StudentClass} \\
    团队成员： &
      \CenteredField[10cm]{\StudentNameA} \\
              &
      \CenteredField[10cm]{\StudentNameB} \\
              &
      \CenteredField[10cm]{\StudentNameC} \\
    指导教师： &
      \CenteredField[10cm]{\Supervisor} \\
    评{\OneEm}{\OneEm}分： &
      \CenteredField[10cm]{\Score} \\
    提交日期： &
      \CenteredField[10cm]{\SubmitDate} \\
  \end{tabular}

  } % 这一行结束 {\Large} 的作用范围
\end{center}


  \vfill
  
  \begin{center}
    {\HeadingFont\fontsize{20pt}{22pt}\selectfont 计算机信息工程学院}
  \end{center}
  
\end{titlepage}

%%%%%%%%%%%%%%%%%%%%
% 前置部分（罗马页码）
%%%%%%%%%%%%%%%%%%%%

\clearpage
\pagenumbering{gobble}

\evenblank
% 学术诚信承诺页
\chapter*{学术诚信承诺}
%\addcontentsline{toc}{chapter}{学术诚信承诺}
\thispagestyle{empty}

{\kai \Large 本人郑重承诺：

本报告及所附代码、图表和数据分析结果，均为本人（或本小组成员）在指导教师指导下独立完成，不含有任何未经注明来源的抄袭、剽窃或其他形式的学术不端行为。

如有违反，本人愿意承担相应的学术与纪律责任。

\vspace{2cm}

\hfill \noindent 完成人（签名）：
\CenteredField[8cm]{}

\vspace{1cm}

\hfill \noindent 日期：
\CenteredField[8cm]{}}

\newpage

% 中文摘要
\chapter*{摘\ 要}
%\addcontentsline{toc}{chapter}{摘\ 要}
\thispagestyle{empty}

本项目完成了教室预约系统的设计与实现。主要包括以下几个要素：

1. \textbf{研究背景与问题概述}：高校教室资源紧张，传统人工预约流程繁琐、信息更新滞后，易出现预约冲突与资源闲置问题，师生难以实时获取教室状态，管理员统筹效率低，亟需智能化系统优化管理。

2. \textbf{数据结构与算法}：采用 “链表 + 栈 + 队列” 组合：链表存储教室信息与预约记录，支持高效增删查；栈记录操作以实现撤销；队列按序处理预约，自动检测冲突。

3. \textbf{主要功能}：分师生与管理员角色，支持教室查询、预约申请与审核、状态实时查看、预约撤销等核心操作。

4. \textbf{实现效果}：核心操作响应高效，支持多用户并发使用，模块化设计便于扩展，显著提升教室资源利用率。

\vspace{1cm}

\noindent\textbf{关键词：} 数据结构；教室预约；链表；栈；队列；C 语言

\newpage

% 英文摘要
\chapter*{Abstract}
%\addcontentsline{toc}{chapter}{Abstract}
\thispagestyle{empty}

The project has completed the design and implementation of a classroom reservation system. It mainly includes the following key elements:

1. \textbf{Background and problem statement}: University classroom resources are in short supply. The traditional manual reservation process is cumbersome with delayed information updates, prone to reservation conflicts and resource idleness. Teachers and students cannot obtain real-time classroom status, and administrators have low coordination efficiency, so an intelligent system is urgently needed to optimize management.

2. \textbf{Data structures and algorithms}: A combination of "linked list + stack + queue" is adopted: linked lists store classroom information and reservation records for efficient addition, deletion and query; stacks record operations for undo; queues process reservations in order with automatic conflict detection

3. \textbf{Main functions}:With teacher/student and administrator roles, it supports core operations such as classroom query, reservation application and review, real-time status check, and reservation cancellation.

4. \textbf{Implementation results}: Core operations respond efficiently, supporting concurrent use by multiple users. The modular design facilitates expansion and significantly improves classroom resource utilization.

Keywords: data structures; classroom reservation; linked list; stack; queue; C language

\vspace{1cm}

\noindent\textbf{Keywords:} data structures; parking lot management; stack; queue; C language

\newpage

% 目录（小写罗马字页码）
\pagenumbering{roman}
\pagestyle{fancy}
\tableofcontents
\thispagestyle{plain}

%%%%%%%%%%%%%%%%%%%%
% 正文部分（阿拉伯数字页码）
%%%%%%%%%%%%%%%%%%%%

\clearpage
\pagestyle{fancy}
\pagenumbering{arabic}
\phantomsection
\setcounter{page}{1}

\cleardoublepage
\chapter{引言}

本章用于向读者介绍本次课程设计的整体情况。建议内容包括：

\begin{itemize}
  \item \textbf{项目背景：} 说明你所选择的应用场景，例如停车场管理、学生成绩管理、图书借阅系统、迷宫求解、表达式求值等。可以简要引用相关领域的背景资料，说明该问题在实际生活或工程中的意义。
  \item \textbf{研究目标：} 用清晰的语句指出本项目希望解决的问题，例如“实现车辆进出的高效管理”“设计支持多种查询操作的学生信息系统”“构建可视化的迷宫寻路算法演示工具”等。
  \item \textbf{核心数据结构：} 简要说明你选择使用的主要数据结构（如栈、队列、链表、树、图等），并用几句话说明选择这些数据结构的理由。
  \item \textbf{报告结构：} 在本节最后，可以用一小段话概括后续各章的内容安排，帮助读者建立整体阅读框架。
\end{itemize}

撰写时请避免把详细技术细节提前堆积在引言中，引言的重点是让读者知道你在干什么、为什么值得做。

\section{项目背景与动机}

这里可以进一步展开应用场景或问题背景。例如：

\begin{itemize}
  \item 城市停车难问题与智能化管理需求；
  \item 校园图书馆借阅系统的数据组织与检索效率；
  \item 表达式计算中栈的应用与重要性；
  \item 图算法在路径规划、社交网络分析中的实际意义等。
\end{itemize}

\section{研究目标与问题定义}

针对本项目，用较为形式化的方式定义任务需求（功能需求、性能需求等），并给出要实现的主要功能模块和预期的性能指标或成功标准。插入图片的方法如图\ref{fig:maodie}。

\begin{figure}[htbp]
  \centering
  \includegraphics[width=0.6\textwidth]{assets/圆头耄耋.png}
  \caption{示例图片}
  \label{fig:maodie}
\end{figure}

下面从“引言”开始的若干页（直到下一章之前）是本次课程设计报告的写作示范（以“引言”章节为例）。请注意，为了模版章节编号不错乱，这篇示例\textbf{没有}显示章节编号。实际写作时仍然应谨遵模版工作。

\chapter*{引言}

随着高校教育规模的不断扩大和教学模式的日益多样化，教室资源的管理与合理分配已成为高校日常运营中的重要课题。传统的教室预约方式通常依赖于人工登记或简单的电子表格，存在信息不透明、预约冲突频发、资源利用率不高等问题。特别是在课程安排灵活、教室类型多样的现代高校环境中，如何实现教室资源的智能化管理和高效利用，成为亟待解决的实际问题。在这样的背景下，本课程设计选取"教室预约系统"作为应用场景，尝试运用数据结构课程中所学的核心知识，设计并实现一个具备基本功能的教室预约管理系统。

本项目的核心目标是利用合适的数据结构，实现教室信息的维护、用户预约、冲突检测、查询统计等基本功能。系统需要支持多种角色的用户（如学生、教师、管理员），能够准确记录教室的使用状态、预约时间、使用者信息等，并根据用户需求智能推荐可用教室。更具体地说，一方面系统要能够高效地处理预约请求，确保时间冲突的自动检测和避免；另一方面，系统需要提供直观的用户交互界面，使不同角色的用户能够方便地进行教室查询、预约、修改和取消操作。项目希望回答的关键问题包括：如何通过合适的数据结构实现教室信息的快速检索；如何设计高效的冲突检测算法以保证预约的准确性；以及如何在保证系统功能完整性的前提下，实现良好的用户体验和系统性能。

在数据结构选择上，本项目的核心设计思想是将教室信息、预约记录、用户信息等分别组织为不同的数据结构。考虑到教室信息需要按多种属性（如教室编号、容量、设备类型等）进行快速检索，系统采用哈希表存储教室基本信息；为支持时间段查询和冲突检测，预约记录采用时间段区间树或按时间排序的链表结构；用户信息则采用平衡二叉树（如AVL树）存储，以保证用户登录和权限验证的效率。这些数据结构的组合充分体现了数据结构课程中关于不同数据结构特性及其适用场景的理解与应用。

本报告的其余部分将按照系统开发的自然逻辑展开：第二章介绍教室预约系统的需求分析，详细说明系统应实现的功能模块及其性能要求；第三章给出系统的总体设计方案，包括数据结构选择、模块划分与接口设计；第四章详细描述各功能模块的算法设计与实现过程，展示关键代码片段并分析其时间复杂度与空间复杂度；第五章展示系统测试结果，通过多组测试用例验证系统功能的正确性和性能的合理性；最后一章对本次课程设计进行总结，反思设计与实现过程中的不足之处，并讨论未来可能的改进方向。
\section*{项目背景与动机}

教室资源是高校教学活动的物质基础，其管理效率直接影响到教学质量与教学秩序。随着高校规模的扩大和教学模式的改革，教室资源的供需矛盾日益突出：一方面，各类课程、社团活动、学术讲座等活动对教室的需求多样化且时间交错；另一方面，教室资源有限，如何实现资源的合理分配和高效利用成为管理难题。传统的人工预约方式存在诸多弊端：预约信息不透明，容易出现"先到先得"的不公平现象；人工处理预约请求效率低下，容易出现遗漏或错误；缺乏有效的冲突检测机制，容易导致预约时间重叠；数据统计困难，难以对教室使用情况进行科学分析。

从计算机科学的角度来看，教室预约管理问题具有典型的数据结构应用特征。教室预约系统本质上是一个资源分配与调度问题，需要处理时间区间管理、冲突检测、多条件查询等复杂操作。这些问题在数据结构课程中都有对应的解决方案：时间区间管理可以使用区间树或线段树；冲突检测可以通过对时间段进行排序和比较来实现；多条件查询可以借助哈希表、平衡二叉树等数据结构。因此，通过教室预约系统的设计与实现，可以将课堂上学习的数据结构知识有机地结合起来，在解决实际问题的过程中加深对理论知识的理解。

选择教室预约系统作为课程设计题目，不仅因为其具有实际应用价值和良好的教学示范作用，还因为该系统具有适度的复杂性，能够较全面地锻炼学生的系统设计能力和编程能力。在功能层面，系统需要实现用户管理、教室信息维护、预约管理、冲突检测、查询统计等多个模块，涉及多种数据结构的综合运用；在性能层面，系统需要保证在高并发情况下的响应速度，特别是在学期初预约高峰期，各项操作的时间复杂度应控制在合理范围内；在工程层面，系统需要具备良好的代码组织结构和清晰的模块划分，便于后续的功能扩展和维护。

\section*{研究目标与问题定义}

本课程设计旨在开发一个高效、可靠的教室预约系统，通过合理的数据结构设计和算法优化，解决高校教室资源分配中存在的实际问题。项目的总体目标包括实现功能完整性、保证系统性能、提升用户体验以及确保数据准确性四个方面。具体而言，系统需要具备教室信息管理、用户预约、冲突检测、审批流程等核心功能，满足师生日常预约需求；通过优化的数据结构选择和算法设计，确保系统在高并发场景下仍能快速响应，关键操作响应时间控制在合理范围内；设计简洁直观的操作界面和流畅的预约流程，降低用户学习成本，提高系统可用性；实现100\%准确的时间冲突检测，避免预约冲突问题，保证教学秩序的正常进行。

针对教室预约系统的实际应用场景，我们定义了需要解决的核心问题。首先是资源分配优化问题，包括如何根据教室的容量、设备配置、位置等因素为不同类型的活动智能推荐合适的教室，以及如何在保证公平性的前提下处理多用户同时预约同一教室的竞争情况。其次是时间冲突检测问题，这涉及到如何高效检测新预约与已有预约的时间冲突、如何支持不同粒度的预约时间单位（如按课时、按小时、按半天等），以及如何处理临时调整和预约变更带来的连锁冲突。第三是数据存储与检索效率问题，需要设计合理的数据存储结构以支持教室信息的快速检索（按编号、容量、位置、设备等多条件查询），管理预约记录的增删改查并在保证操作效率的同时维持数据一致性，还要支持历史数据的统计分析，为教室资源优化提供数据支持。最后是多角色权限管理问题，需要设计合理的权限体系来区分学生、教师、管理员等不同角色的操作权限，实现灵活的审批流程以支持不同场景下的审核要求，并记录操作日志以确保系统的可追溯性和安全性。

为明确项目边界，我们界定了研究范围和限制。功能范围包括支持教室基本信息管理（添加、修改、查询、删除）、用户注册登录和权限管理、预约申请审批修改取消全流程、时间冲突自动检测、预约记录查询和统计分析以及系统基本设置和管理。技术限制方面，采用C语言实现且不依赖复杂的第三方库，使用SQLite作为数据存储以保证系统的轻量级和易部署性，暂不考虑网络分布式部署而以单机应用为主要设计目标，暂不提供图形用户界面而通过命令行界面进行功能演示。性能目标则要求支持至少100个教室信息的存储和管理、至少500个用户的同时使用、至少1000条预约记录的存储和查询，同时确保关键操作响应时间符合性能需求分析中定义的标准。

\section*{核心数据结构选择}

在教室预约系统的设计中，我们根据数据的不同特点和操作需求，精心选择了一系列数据结构，使每种结构都能在最适合的场景下发挥最大效能。这就像为不同的工作任务选择合适的工具一样，既保证了系统的运行效率，又提升了代码的可维护性。

对于教室基本信息，我们采用了哈希表作为主要的存储结构。这是因为教室编号通常作为查询教室信息的关键依据，而哈希表能够以平均O(1)的时间复杂度完成基于键值的快速查找，非常符合这一需求。当用户输入教室编号时，系统能够几乎立即返回对应的教室详情，包括容量、位置、设备配置等信息。此外，为支持按其他属性（如容量范围、位置区域等）的查询，我们还建立了相应的辅助索引结构，确保多维度的查询需求都能得到高效满足。

对于预约记录的管理，时间冲突检测是最核心的操作。我们选择使用区间树（Interval Tree）来存储每个教室的预约时间段。这种数据结构专门用于处理区间查询问题，能够在O(log n)的时间内判断新预约时间段是否与已有预约重叠，从而快速、准确地完成冲突检测。如果考虑到实现复杂度，按开始时间排序的有序链表结合二分查找也是一种可行的备选方案，同样能保证较高的查询效率。

用户信息需要频繁地进行查找、插入和更新操作，我们选择了平衡二叉树（AVL树）来存储这些数据。AVL树能够始终保持树的平衡，确保在最坏情况下查找、插入和删除操作的时间复杂度均为O(log n)。这对于用户登录验证、信息修改等场景非常有利，即使系统中的用户数量增长到数千甚至上万，系统仍能保持稳定的响应速度。

除此之外，系统还使用了队列来管理待审批的预约请求。这种“先进先出”的结构确保了所有请求都能按照提交顺序得到公平处理，避免了优先级混乱的问题。我们还用栈来记录用户的操作历史，当用户误操作时，可以通过撤销功能回退到之前的状态，提升了系统的容错性和用户体验。最后，图结构被用于表示教室之间的位置关系，当用户需要查找“最近的可用教室”时，系统能够基于图算法提供智能推荐，使资源分配更加合理。

通过这种“哈希表+区间树+AVL树”的组合，并辅以队列、栈、图等辅助结构，我们为教室预约系统构建了一个既高效又灵活的数据层。每种数据结构都各司其职，协同工作，共同支撑起了系统的各项功能，为后续的算法实现和性能优化奠定了坚实的基础。


\chapter{需求分析与系统设计}

本章重点在于把系统需求讲清楚，并给出总体设计方案。读者在看完本章后，应当能够理解系统应该具备哪些功能，以及采用怎样的整体架构来实现这些功能。

\section{功能需求分析}

教室预约系统需要为三类用户提供相应的功能支持：学生、教师和管理员。学生用户可以查询可用教室并提交预约申请，教师拥有优先预约权限并可能参与审批流程，管理员则负责教室信息维护、系统监控和统计分析等管理职能。

从功能模块来看，系统主要包括四个核心部分：教室信息管理（维护教室静态信息及可用状态）、预约流程管理（覆盖申请、审批、使用等完整生命周期）、冲突检测与处理（确保时间冲突的准确识别和避免）、查询统计功能（支持多维度数据检索和分析）。此外，系统还应提供预约提醒、操作历史记录、消息通知等辅助功能，以提升用户体验和系统实用性。

总体而言，系统设计应以用户需求为中心，在确保核心功能稳定可靠的前提下，提供简洁直观的操作流程和必要的智能化辅助，实现高效、公平、便捷的教室资源管理。

\section{性能需求分析}

为确保系统稳定高效运行，需从多个维度设定明确的性能指标。响应速度方面，常规操作应在2秒内完成，复杂查询不超过5秒，冲突检测在1秒内完成，页面加载时间控制在3秒以内。并发处理能力要求支持至少100名用户同时在线，在高负荷时段仍能保持稳定运行。

数据准确性方面，时间冲突检测必须达到100\%准确，核心数据必须保持一致性，关键操作需采用事务处理机制。系统将采用SQLite数据库，利用其轻量级和事务支持特性确保数据安全。系统可靠性要求达到99%以上的正常服务时间，具备良好的容错能力和数据备份恢复机制。

可扩展性要求系统采用模块化设计，支持新增属性和功能扩展，并提供开放接口以便与其他系统集成。用户体验方面，界面应直观友好，流程简洁明了，操作反馈及时明确，并提供完整的帮助文档。

通过以上性能指标，我们旨在构建一个响应迅速、运行稳定、数据准确、易于使用且具备良好扩展性的教室预约系统，切实满足高校师生的实际需求。

\section{总体设计方案}

系统采用模块化设计思想，主要包括以下几个功能模块：

\begin{itemize}
  \item \textbf{数据结构模块：} 实现栈、队列、链表等基础数据结构；
  \item \textbf{车辆管理模块：} 处理车辆的到达、离开、查询等操作；
  \item \textbf{费用计算模块：} 根据停车时长计算费用；
  \item \textbf{显示统计模块：} 显示停车场和便道的当前状态；
  \item \textbf{主控模块：} 提供用户交互界面，调用各功能模块完成任务。
\end{itemize}

下面给出系统功能模块图的示例（使用 \LaTeX 绘图功能）：

\begin{figure}[htbp]
  \centering
  \begin{tikzpicture}[
    node distance=1cm and 1.2cm,
    module/.style={
      rectangle,
      rounded corners,
      draw=blue!60,
      fill=blue!5,
      thick,
      minimum width=3cm,
      minimum height=1cm,
      text centered,
      font=\small
    },
    submodule/.style={
      rectangle,
      rounded corners,
      draw=green!60,
      fill=green!5,
      thick,
      minimum width=1.2cm,
      minimum height=0.8cm,
      text centered,
      font=\small
    },
    mainmodule/.style={
      rectangle,
      rounded corners,
      draw=red!60,
      fill=red!10,
      very thick,
      minimum width=3.5cm,
      minimum height=1.2cm,
      text centered,
      font=\small\bfseries
    },
    arrow/.style={
      -Stealth,
      thick,
      draw=gray!70,
      rounded corners=3pt
    }
  ]

  % 主控模块（最上层）
  \node[mainmodule] (main) at (0, 0) {主控模块（界面）};

  % 第二层：功能模块
  \node[module, below left=of main, xshift=-1.5cm] (vehicle) {车辆管理模块};
  \node[module, below=2cm of main] (fee) {费用计算模块};
  \node[module, below right=of main, xshift=1.5cm] (display) {显示统计模块};

  % 第三层：车辆管理子模块
  \node[submodule, below=of vehicle, xshift=-1.5cm] (arrive) {到达};
  \node[submodule, below=of vehicle, xshift=0cm] (depart) {离开};
  \node[submodule, below=of vehicle, xshift=1.5cm] (query) {查询};

  % 第三层：显示统计子模块
  \node[submodule, below=of display, xshift=-1cm] (status) {状态显示};
  \node[submodule, below=of display, xshift=1cm] (report) {报表生成};

  % 底层：数据结构模块
  \node[module, below=5cm of main] (datastructure) {数据结构模块};

  \node[submodule, below=of datastructure, xshift=-3cm] (stack) {停车场栈};
  \node[submodule, below=of datastructure, xshift=0cm] (queue) {便道队列};
  \node[submodule, below=of datastructure, xshift=3cm] (list) {车辆链表};

  % ===== 定义中间转折高度 =====
  \coordinate (midlevel) at (0, -5);  % 定义统一的中间转折线高度

  % ===== 智能肘线连接 =====
  
  % 主控到功能模块
  \draw[arrow] (main.south) |- (0, -1) - | (vehicle.north);
  \draw[arrow] (main) -- (fee);  % 垂直对齐，用直线
  \draw[arrow] (main.south) |- (0, -1) - | (display.north);

  % 车辆管理到子模块
  \draw[arrow] (vehicle.south) |- (-6.3, -3.1) - | (arrive.north);
  \draw[arrow] (vehicle) -- (depart);  % 垂直对齐
  \draw[arrow] (vehicle.south) |- (-5.5, -3.1) - | (query.north);

  % 显示统计到子模块
  \draw[arrow] (display.south) |- (5.5, -3.1) - | (status.north);
  \draw[arrow] (display.south) |- (6.3, -3.1) - | (report.north);

  % 功能模块到数据结构
  % 方法：起点 |- 中间高度与目标x坐标的交点 -- 终点
  \draw[arrow] (arrive.south) |- (midlevel -| datastructure.north) -- (datastructure.north);
  \draw[arrow] (depart.south) |- (midlevel -| datastructure.north) -- (datastructure.north);
  \draw[arrow] (query.south) |- (midlevel -| datastructure.north) -- (datastructure.north);
  \draw[arrow] (status.south) |- (midlevel -| datastructure.north) -- (datastructure.north);
  
  % 费用计算直接连接（垂直对齐）
  \draw[arrow] (fee) -- (datastructure);

  % 数据结构到具体实现
  \draw[arrow] (datastructure.south) |- (0, -7.1) - | (stack.north);
  \draw[arrow] (datastructure) -- (queue);  % 垂直对齐
  \draw[arrow] (datastructure.south) |- (0, -7.1) - | (list.north);

  % 费用计算与车辆离开的连接
  \coordinate (feedep) at ($ (depart.north) + (0, 0.2) $);
  \draw[-Stealth, thick, dashed, draw=orange, rounded corners=3pt] 
    (depart.north) -- (feedep) |- (fee.west);

  % 添加背景框
  \begin{scope}[on background layer]
    \node[
      rectangle,
      rounded corners,
      draw=gray!30,
      fill=gray!2,
      fit=(vehicle)(fee)(display)(arrive)(depart)(query)(status)(report),
      inner sep=0.3cm,
      label={[anchor=north west, xshift=1.5cm, yshift=-0.2cm]above:业务逻辑层}
    ] {};

    \node[
      rectangle,
      rounded corners,
      draw=gray!30,
      fill=gray!2,
      fit=(datastructure)(stack)(queue)(list),
      inner sep=0.3cm,
      label={[anchor=north west, xshift=2cm, yshift=-0.2cm]above:数据层}
    ] {};
  \end{scope}

  \end{tikzpicture}
  \caption{停车场管理系统功能模块图}
\end{figure}


\section{数据结构设计}

\subsection{教室预约栈的设计}

教室预约采用顺序栈实现，主要结构定义如下：

\begin{codeblock}
// 车辆信息结构
typedef struct {
    char plate[20];     // 车牌号
    time_t arriveTime;  // 到达时间
    int position;       // 停车位置
} Car;

// 停车场栈结构
typedef struct {
    Car data[MAX_SIZE]; // 栈数组
    int top;            // 栈顶指针
    int capacity;       // 栈容量
} ParkingStack;
\end{codeblock}

\subsection{便道队列的设计}

便道采用链式队列实现，主要结构定义如下：

\begin{codeblock}
// 队列节点
typedef struct QNode {
    Car data;
    struct QNode *next;
} QueueNode;

// 链式队列
typedef struct {
    QueueNode *front;  // 队头指针
    QueueNode *rear;   // 队尾指针
    int size;          // 队列长度
} WaitingQueue;
\end{codeblock}

\cleardoublepage
\chapter{详细设计与算法实现}

本章详细描述各功能模块的算法设计思路，展示关键代码实现，并对算法的时间复杂度和空间复杂度进行分析。

\section{栈和队列的基本操作}

\subsection{栈的初始化与基本操作}

栈的初始化非常简单，只需将栈顶指针设为 -1 即可：

\begin{codeblock}
// 初始化停车场栈
void initStack(ParkingStack *s, int capacity) {
    s->top = -1;
    s->capacity = capacity;
}

// 判断栈是否为空
int isEmpty(ParkingStack *s) {
    return s->top == -1;
}

// 判断栈是否已满
int isFull(ParkingStack *s) {
    return s->top == s->capacity - 1;
}
\end{codeblock}

压栈操作将新元素添加到栈顶，时间复杂度为 $O(1)$：

\begin{codeblock}
// 车辆进入停车场（压栈）
int push(ParkingStack *s, Car car) {
    if (isFull(s)) {
        return 0;  // 停车场已满
    }
    s->data[++s->top] = car;
    s->data[s->top].position = s->top + 1;
    return 1;
}
\end{codeblock}

\subsection{队列的初始化与基本操作}

链式队列的初始化需要创建头节点：

\begin{codeblock}
// 初始化便道队列
void initQueue(WaitingQueue *q) {
    q->front = q->rear = NULL;
    q->size = 0;
}

// 入队操作
void enqueue(WaitingQueue *q, Car car) {
    QueueNode *node = (QueueNode*)malloc(sizeof(QueueNode));
    node->data = car;
    node->next = NULL;
    
    if (q->rear == NULL) {
        q->front = q->rear = node;
    } else {
        q->rear->next = node;
        q->rear = node;
    }
    q->size++;
}
\end{codeblock}

\section{车辆到达处理算法}

车辆到达时的处理流程如下：

\begin{enumerate}
  \item 检查停车场是否已满；
  \item 如果未满，将车辆信息压入停车场栈；
  \item 如果已满，将车辆信息加入便道队列。
\end{enumerate}

算法的伪代码表示：

\begin{codeblock}
// 车辆到达处理
void carArrival(ParkingStack *parking, 
                WaitingQueue *waiting, 
                char *plate) {
    Car car;
    strcpy(car.plate, plate);
    car.arriveTime = time(NULL);
    
    if (!isFull(parking)) {
        push(parking, car);
        printf("车辆 %s 已进入停车场，车位号 %d\n", 
               plate, parking->top + 1);
    } else {
        enqueue(waiting, car);
        printf("停车场已满，车辆 %s 进入便道等待\n", 
               plate);
    }
}
\end{codeblock}

\textbf{复杂度分析：} 该算法的时间复杂度为 $O(1)$，空间复杂度为 $O(1)$（不考虑数据结构本身占用的空间）。

\section{车辆离开处理算法}

车辆离开处理是系统中最复杂的操作，需要考虑以下情况：

\begin{enumerate}
  \item 查找目标车辆在停车场中的位置；
  \item 如果车辆不在栈顶，需要将其上方的车辆暂时移出；
  \item 目标车辆离开后，将暂时移出的车辆依次放回；
  \item 如果便道有车等待，让队首车辆进入停车场。
\end{enumerate}

关键代码如下：

\begin{codeblock}
// 车辆离开处理
void carDeparture(ParkingStack *parking, 
                  WaitingQueue *waiting, 
                  char *plate) {
    ParkingStack temp;  // 临时栈
    initStack(&temp, parking->capacity);
    
    int found = 0;
    Car leavingCar;
    
    // 查找目标车辆
    while (!isEmpty(parking)) {
        Car car = pop(parking);
        if (strcmp(car.plate, plate) == 0) {
            leavingCar = car;
            found = 1;
            break;
        }
        push(&temp, car);  // 暂存到临时栈
    }
    
    // 将临时栈的车辆放回停车场
    while (!isEmpty(&temp)) {
        push(parking, pop(&temp));
    }
    
    if (found) {
        // 计算停车费用
        double fee = calculateFee(leavingCar.arriveTime);
        printf("车辆 %s 离开，停车费用 %.2f 元\n", 
               plate, fee);
        
        // 如果便道有车，让队首进入
        if (waiting->size > 0) {
            Car waitingCar = dequeue(waiting);
            push(parking, waitingCar);
            printf("便道车辆 %s 进入停车场\n", 
                   waitingCar.plate);
        }
    }
}
\end{codeblock}

\textbf{复杂度分析：} 在最坏情况下（目标车辆在栈底），需要遍历整个停车场栈，时间复杂度为 $O(n)$，其中 $n$ 为停车场容量。空间复杂度为 $O(n)$，因为需要使用临时栈。

\cleardoublepage
\chapter{测试与分析}

本章展示系统的测试过程和测试结果，通过多组测试用例验证系统功能的正确性。

\section{测试环境}

\begin{itemize}
  \item \textbf{操作系统：} Windows 11 / Ubuntu 22.04
  \item \textbf{编译器：} GCC 11.4.0 / MinGW-w64 GCC 13.1.0
  \item \textbf{开发工具：} Visual Studio Code / 编译方式：XeLaTeX
  \item \textbf{编译选项：} \texttt{gcc -o parking parking.c -std=c11 -Wall}
\end{itemize}

\section{测试用例设计}

本项目设计了以下几组测试用例来验证系统功能，详见表\ref{tab:testcase}。

\begin{table}[htbp]
  \centering
  \caption{测试用例说明}
  \label{tab:testcase}
  \begin{tabular}{lll}
    \toprule
    \textbf{测试编号} & \textbf{测试目的} & \textbf{测试内容} \\
    \midrule
    TC-01 & 基本进出功能 & 车辆正常进入和离开停车场 \\
    TC-02 & 停车场满载测试 & 停车场满时车辆进入便道 \\
    TC-03 & 便道车辆进入 & 停车场有空位时便道车辆自动进入 \\
    TC-04 & 车辆查询功能 & 按车牌号查询车辆信息 \\
    TC-05 & 边界条件测试 & 空停车场、满停车场等边界情况 \\
    TC-06 & 费用计算测试 & 验证不同停车时长的费用计算 \\
    \bottomrule
  \end{tabular}
\end{table}

\section{测试结果展示}

以下是一组典型测试用例的运行结果：

\begin{codeblock}
========== 停车场管理系统 ==========
停车场容量：5

[操作] 车辆 A001 到达
车辆 A001 已进入停车场，车位号 1

[操作] 车辆 A002 到达
车辆 A002 已进入停车场，车位号 2

[查询] 停车场状态
停车场：5/2 (总数/已用)
便道等待：0 辆

[操作] 车辆 A001 离开
车辆 A001 离开，停车费用 15.00 元
\end{codeblock}

\section{性能分析}

表\ref{tab:complex}对系统各主要操作的时间复杂度进行总结。

\begin{table}[htbp]
  \centering
  \caption{算法时间复杂度分析}
  \label{tab:complex}
  \begin{tabular}{lcc}
    \toprule
    \textbf{操作} & \textbf{最好情况} & \textbf{最坏情况} \\
    \midrule
    车辆到达 & $O(1)$ & $O(1)$ \\
    车辆离开（在栈顶） & $O(1)$ & $O(1)$ \\
    车辆离开（在栈底） & $O(n)$ & $O(n)$ \\
    车辆查询 & $O(1)$ & $O(n)$ \\
    状态显示 & $O(1)$ & $O(1)$ \\
    \bottomrule
  \end{tabular}
\end{table}

\cleardoublepage
\chapter{总结与展望}

\section{工作总结}

本课程设计完成了一个基于栈、队列、链表等数据结构的教室预约管理系统。系统采用模块化设计思想，将停车场栈和便道队列的基本操作封装为独立模块，实现了教室信息的维护、用户预约、冲突检测、查询统计等核心功能模块。通过该项目的设计与实现，将课堂上学习的数据结构理论知识应用到了实际问题的解决中，加深了对栈、队列、链表等基本数据结构的理解，同时也锻炼了C语言编程能力和系统设计能力。

在数据结构方面，通过合理选择哈希表、区间树和AVL树等数据结构，充分利用了它们的特性来模拟教室预约的实际场景。哈希表的快速查找特性完美契合了按教室编号查询的需求，区间树的高效区间查询很好地解决了时间冲突检测问题，AVL树的平衡特性则保证了用户信息管理的效率。在算法设计方面，实现了用户注册登录、教室预约、冲突检测、审批流程等核心功能，并对关键算法的时间复杂度和空间复杂度进行了分析，确保了系统在性能上满足实际应用的基本要求。

\section{主要收获}

通过本次课程设计，主要有以下几方面的收获：

\begin{enumerate}
  \item \textbf{数据结构应用能力：} 学会了如何根据实际问题的特点选择合适的数据结构，并将理论知识应用到具体的系统设计中。特别是在教室预约系统中，深刻理解了哈希表、区间树、平衡二叉树等数据结构的适用场景和优缺点。
  \item \textbf{编程实践能力：} 熟练掌握了C语言中结构体、指针、文件操作等重要概念的使用，提高了代码编写和调试能力。通过SQLite数据库的集成使用，掌握了数据库的基本操作和数据持久化方法。
  \item \textbf{系统设计思维：} 学会了将复杂问题分解为多个功能模块，通过模块化设计提高代码的可读性和可维护性。特别是在处理多角色权限管理、预约审批流程等复杂业务逻辑时，学会了如何设计清晰的系统架构。
  \item \textbf{算法分析能力：} 掌握了如何分析算法的时间复杂度和空间复杂度，能够评估不同算法方案的优劣。特别是在冲突检测算法的设计中，深入理解了区间查询算法的效率差异。
  \item \textbf{问题解决能力：} 在项目实施过程中遇到了许多实际问题（如边界条件处理、内存管理、并发控制等），通过查阅资料和实践探索最终得到了解决。特别是在时间冲突检测的准确性保证上，进行了多次算法优化和测试验证。
\end{enumerate}

\section{不足与改进方向}

虽然本系统实现了基本的教室预约管理功能，但仍存在一些不足之处，未来可从以下几个方面进行改进：

\begin{enumerate}
  \item \textbf{界面优化：} 当前系统采用命令行交互方式，用户体验有限。可以考虑开发图形化用户界面，提供更加直观友好的操作方式，特别是对于非技术背景的用户群体。
  \item \textbf{功能扩展：} 系统目前主要关注教室预约的基本功能，可以增加教室使用统计分析、预约趋势预测、智能推荐等功能，为教室资源优化配置提供数据支持。
  \item \textbf{性能优化：} 在数据量较大时，部分查询操作的时间复杂度较高。可以引入更高效的数据结构（如B+树索引）或缓存机制，提升系统的响应速度和处理能力。
  \item \textbf{移动端支持：} 随着移动设备的普及，可以开发移动应用程序，方便用户随时随地查询教室信息和提交预约申请，提高系统的便捷性和可用性。
  \item \textbf{系统集成：} 目前系统独立运行，未来可以考虑与学校现有的教务系统、身份认证系统等进行集成，实现单点登录和数据共享，减少用户重复输入信息的麻烦。
  \item \textbf{安全性增强：} 虽然系统采用了基本的权限控制，但在数据加密、防止SQL注入、防篡改等方面还可以进一步加强，提高系统的安全性。
\end{enumerate}

总的来说，通过本次课程设计的学习和实践，不仅巩固了数据结构的理论知识，也培养了实际编程能力和工程实践能力，为今后进一步学习更复杂的数据结构和算法打下了良好的基础。
\clearpage
\cleardoublepage

% 以下是参考文献部分
\phantomsection
\addcontentsline{toc}{chapter}{参考文献}

% {99} 是最宽的编号，为两位数编号预留空间
\begin{thebibliography}{99}

% 类型：专著 [M]
\bibitem[1]{yanweimin} 严蔚敏, 李冬梅, 吴伟民. 数据结构（C语言版）[M]. 第2版. 北京: 人民邮电出版社, 2015.

% 类型：专著 [M]
\bibitem[2]{wanghongmei} 王红梅, 胡明, 王涛. 数据结构（C语言版）学习指导[M]. 北京: 清华大学出版社, 2019.

% 类型：专著 [M]
\bibitem[3]{yinrenkun} 殷人昆. 数据结构（用面向对象方法与C++语言描述）[M]. 第2版. 北京: 清华大学出版社, 2007.

% 类型：联机资源 [EB/OL]
\bibitem[4]{cprog} C语言中文网. C语言教程[EB/OL]. [引用日期 2025-12-20]. 可获取自: \url{http://c.biancheng.net/}

\end{thebibliography}

\appendix

\cleardoublepage
\chapter{附录说明}

附录部分可以包含：

\begin{itemize}
  \item 完整的源代码清单（主体报告中只保留关键片段）；
  \item 详细的测试数据和运行日志；
  \item 系统流程图、数据流图等补充图表；
  \item 开发过程中遇到的问题及解决方案记录。
\end{itemize}

\section{主要源代码}

这里可以附上主要的头文件和源文件：

\begin{codeblock}
// parking.h - 头文件
#ifndef PARKING_H
#define PARKING_H

#include <stdio.h>
#include <stdlib.h>
#include <string.h>
#include <time.h>

#define MAX_SIZE 100
#define MAX_PLATE 20

// 车辆信息结构
typedef struct {
    char plate[MAX_PLATE];
    time_t arriveTime;
    int position;
} Car;

// 停车场栈结构
typedef struct {
    Car data[MAX_SIZE];
    int top;
    int capacity;
} ParkingStack;

// 队列节点
typedef struct QNode {
    Car data;
    struct QNode *next;
} QueueNode;

// 便道队列
typedef struct {
    QueueNode *front;
    QueueNode *rear;
    int size;
} WaitingQueue;

// 函数声明
void initStack(ParkingStack *s, int capacity);
int push(ParkingStack *s, Car car);
Car pop(ParkingStack *s);
void initQueue(WaitingQueue *q);
void enqueue(WaitingQueue *q, Car car);
Car dequeue(WaitingQueue *q);

#endif
\end{codeblock}

\section{测试数据示例}

以下是部分测试数据及预期输出：

\begin{codeblock}
测试场景：停车场容量为3，依次有5辆车到达

输入序列：
A A001
A A002
A A003
A A004
A A005

预期输出：
车辆 A001 已进入停车场，车位号 1
车辆 A002 已进入停车场，车位号 2
车辆 A003 已进入停车场，车位号 3
停车场已满，车辆 A004 进入便道等待
停车场已满，车辆 A005 进入便道等待
\end{codeblock}

\end{document}
